\section{Discussion}
\label{sec:discussion}

\paragraph{What the evidence says.} The gains that learning-guided search reports are mostly \emph{structure} (destruction--reconstruction), \emph{budget} and \emph{tuning}. The learner's own contribution, the decision effect, is bounded by one quantity, the headroom of the best fixed action net of run-to-run noise. Below about $0.05$ points the online learner never gained beyond the equivalence margin in any cell outside those that set the threshold. Above it, it sometimes recovers part of what a validated action gives and often loses. Even where the headroom is large (unrelated machines) a learner captures a small fraction of it with uniform times. The apparent exception, the standard Q-Learning with correlated machines (up to about two thirds of the validated gain), disappears when the order of the actions is randomised: the standard learner repeats the action listed first, so its gains and losses measure where a good action stood in the list (Section~\ref{sec:order}). The published designs and the irace-tuned learner do not change the order. What headroom rewards is knowledge acquired \emph{before} the run: a validated action, or a selector trained offline, which captures about four fifths of the one-seed oracle gain on $P||\Cmax$ and $60$--$73\%$ on unrelated machines, mostly through class-level knowledge. Where the actions are complementary (the genetic algorithm, the ILS operators) the static competitor must be a \emph{mixture}, not a single action; the claim ``a validated static configuration is never beaten beyond the margin'' is a claim about the best of a single action or a validated mixture.

\paragraph{Why the usual comparison misleads.} The comparison of a learning-guided metaheuristic with a non-learning one changes at least four things at once: the search structure, the number of iterations or the CPU time, the tuning effort that went into each side, and the information the learner receives. The protocol varies one at a time. Here the learning-guided IG appeared to owe its advantage over the ILS to the learner, when it belongs to the destruction--reconstruction mechanism; at equal CPU time the learner appeared to be responsible for a gain that belongs to choosing cheaper actions; and a one-seed oracle appeared to show room for improvement that was mostly noise. Omitting any one of these controls leaves room for a positive result that is not the learner's. In the literature we audited ($23$ recent open-access papers) $22$ report a benefit of the learned component, $10$ include a random-selection control, $6$ compare at a common time limit and $3$ include a static configuration. We do not claim that those benefits are absent: the audit counts which controls were reported, not whether the claims are true.

\paragraph{A rule of use.} Estimate the cross-fitted headroom on tuning instances (Section~\ref{sec:estimator}) or its cheap proxy, the validated-gain-to-noise ratio (AUC $0.75$ and $0.73$ for ``Q-Learning gains'' on the $26$ (subset, $m$, $n$) classes of an unseen family, F). Below about $0.05$ points, Q-Learning does not pay. In all cases a validated action or mixture, or better a policy trained offline, is the baseline to beat. The rule is a screening device, not a law: the threshold was set on the discovery cells, it does not discriminate within the flowshop, at finer times or on unrelated machines, and the base-rate, bootstrap and one-sided readings of it were post hoc.

\paragraph{A checklist for a claim that a learned selector helps.}
\begin{enumerate}
\item Report a \textbf{random control over the same actions at the same budget}, with the same code path as the learner.
\item At equal CPU time, \textbf{calibrate the cost ratio} $K$ on separate instances, report it for every arm, and split the gain into a decision effect and a cost effect.
\item \textbf{Validate a static action (or a mixture) offline} on instances disjoint from the test set, per class, with the same effort the learner's tuning received; add a configurator (irace) and a published AOS design as competitors where they apply.
\item Estimate the headroom \textbf{with cross-fitting}, and report both the one-seed value (an upper bound) and the cross-fitted value (a lower bound).
\item Include a \textbf{positive control} (a planted effect the protocol must detect) and an \textbf{equivalence margin}, so that ``no effect'' is a finding and not a failure to reject.
\item \textbf{Record hypotheses, decision rules and predictions before running}, with hashes of the code and frozen inputs, and control the error rate within each family of tests.
\end{enumerate}

\paragraph{What the registered predictions showed.} The batteries were specified before they were run; \npredheld{} of \npredtotal{} point predictions held in \nbatteries{} batteries. The misses (cost ratios, the classic AOS methods, the flowshop, warm-started Q-Learning, finer times) led to the headroom rule and to cross-fitting. The headline pairs keep their significance under Pratt's treatment of ties and a permutation test, and $663$ of the $685$ confirmed tests survive a global Benjamini--Yekutieli control over the $1{,}169$ tests of all registered batteries.

\section{Limitations}
\label{sec:limitations}
\begin{enumerate}[label=(\roman*)]
\item \emph{Equal CPU time.} The comparison depends on $K$, measured on separate instances in sequence (not always on a quiet machine) and reported for every arm. A learner whose actions are cheap can win on cost alone, which is why the decomposition is part of the protocol.
\item \emph{Practical significance.} Gains are in points of the makespan, so $0.05\%$ may or may not matter, and the margin is a choice. As an absolute reference, the IG with random choice reaches the classical lower bound in $42$--$48\%$ of the synthetic runs; we did not compare it with the solvers reviewed in~\cite{OstojicEtAl2025}, because the protocol compares arms, not absolute quality.
\item \emph{Registration and multiplicity.} Each battery has its own Holm family and registration; the headroom rule began as an exploratory reading and rests on few cells per space and host. The pre-specification is a dated, hash-stamped diary under version control, not a public registry (Appendix~\ref{app:batteries}).
\item \emph{Scope.} The evidence is for $P||\Cmax$, $R||\Cmax$ (synthetic, with three published time structures and with the standard benchmark), the permutation flowshop and one genetic algorithm, all with a short list of tabular and bandit learners, a deep Q-network on joint actions, and two published AOS designs. Deep policy gradients, policies with a learned state trained across instances with a neural network, and portfolios of hundreds of operators are outside the evidence; the pre-trained Q table of Section~\ref{sec:repairs} is the tabular form of such a policy and did not help.
\item \emph{Horizons.} Horizons up to $T=10^5$ were studied on two hosts; a learner that needs a much longer horizon than the headroom lasts for is not excluded in general, but it would have to be shown to pay at a budget where the headroom is already gone.
\item \emph{Estimators.} The cross-fitted headroom is a conservative, downward-biased estimate (Proposition~\ref{prop:bounds}); the one-seed oracle is an upper bound. Neither is the true headroom, and the rule of use is stated for the cross-fitted one.
\item \emph{Selector cost.} The offline selector's training cost is not charged to it in the main comparisons ($85{,}800$ and $257{,}400$ IG runs for unrelated machines, about $0.7$ CPU-hours); amortised, it amounts to $21{,}700$ and $65{,}200$ extra iterations per held-out instance.
\item \emph{Review status.} The manuscript has not been externally reviewed; simulated blind reviews guided the revisions but are not a substitute.
\end{enumerate}

\section{Conclusion}
\label{sec:conclusion}
A learned operator selector is not what explains the advantage of an IG over ILS for identical parallel machines: that belongs to destruction--reconstruction, and the share of a local search's gain that is only budget ($93\%$, $30\%$ or $28\%$ by implementation) makes the CPU cost ratio a result to report. The oracle headroom net of run-to-run noise is a necessary condition for Q-Learning to pay, not a sufficient one: no cell below about $0.05$ points gains beyond the margin, while above it Q-Learning gains $0.15$--$0.25$ points in the middle and upper size terciles of family C (confirmed on held-out data) and loses in the small classes. A validated static configuration was never worse than the online learner beyond the margin, even on unrelated machines and on published instance families, and a selector trained offline exploits the headroom, mostly through class-level knowledge; four repairs of the online learner do not help, and its limit is its feedback.

For the practitioner the message is a procedure more than a verdict: before crediting a learned component, compare it with random choice at the same budget, with a static choice validated offline and with a configurator, and ask whether the headroom is large enough for any online method to find. For the researcher it is a set of estimators (the cost decomposition and the cross-fitted headroom) whose statistical properties are stated and checked, and a registry of $\npredtotal$ pre-recorded predictions, whose misses are listed and changed the analysis.

\paragraph{Data and code availability.} Code, instances, results, the research diary with the pre-registered hypotheses and md5 hashes, the protocol library and the scripts that recompute every number are at \url{https://github.com/CarlosAraujo7/pmsp}. \todonote{funding statement: CAPES/FUNCAP/CNPq to confirm; anonymised repository link for review}
